\chapter{Principio de inclusión-exclusión}
\label{md:inclusion-exclusion}

\section{Uniones y condiciones prohibidas}

Para conjuntos finitos $A,B$,
\[
 |A\cup B|=|A|+|B|-|A\cap B|.
\]
Se resta la intersección porque la suma inicial la cuenta dos veces.
Para $A_1,\ldots,A_r$, el \textbf{principio de inclusión-exclusión} es
\[
 \boxed{\left|\bigcup_{i=1}^r A_i\right|
 =\sum_{k=1}^r(-1)^{k-1}
   \sum_{1\le i_1<\cdots<i_k\le r}
   |A_{i_1}\cap\cdots\cap A_{i_k}|.}
\]
Se suman los conjuntos individuales, se restan las intersecciones
de dos, se suman las de tres, y se continúa alternando signos.
Cada grupo de índices se cuenta una sola vez.

Si $A_i\subseteq S$, con $S$ finito, representa el cumplimiento
de una condición $i$, el número de elementos que no cumplen ninguna es
\[
 \left|S\setminus\bigcup_{i=1}^r A_i\right|
 =|S|+\sum_{k=1}^r(-1)^k
   \sum_{1\le i_1<\cdots<i_k\le r}
   |A_{i_1}\cap\cdots\cap A_{i_k}|.
\]
Es frecuente que las $A_i$ sean condiciones prohibidas:
sus complementos simultáneos constituyen los resultados válidos.

La idea de la prueba es seguir la contribución de un elemento que
pertenezca a exactamente $t$ conjuntos. En la fórmula del complemento
se cuenta $\sum_{k=0}^t(-1)^k\binom tk=(1-1)^t$ veces:
una vez si $t=0$, y cero si $t\ge1$.

Para aplicar el principio se fija el universo, se definen las
condiciones y se cuentan sus intersecciones. No basta conocer solo
los tamaños individuales, salvo que se sepa que no hay solapamientos.

\begin{ejemplo}
Entre $1$ y $30$, hay $15$ múltiplos de $2$, $6$ de $5$ y $3$ de ambos.
Los divisibles por alguno de los dos son $15+6-3=18$;
los que no son divisibles por ninguno son $30-18=12$.
\end{ejemplo}

\newpage
\section{Desordenaciones}
\label{md:desordenaciones}

Una \textbf{desordenación} de $n$ elementos es una permutación sin
puntos fijos: $\sigma(i)\ne i$ para todos los índices.
Sea $d(n)$ su número.
En el universo de las $n!$ permutaciones, la condición prohibida
$A_i$ es dejar fijo $i$. Fijar $k$ índices deja $(n-k)!$ permutaciones,
y hay $\binom nk$ maneras de elegir esos índices.
Por inclusión-exclusión,
\[
 \boxed{d(n)=\sum_{k=0}^n(-1)^k\binom nk(n-k)!
          =n!\sum_{k=0}^n\frac{(-1)^k}{k!}.}
\]
Las permutaciones con al menos un punto fijo son $n!-d(n)$.
El recuento tiene una relación recurrente equivalente, presentada
en la sección~\ref{md:modelar-recurrencia}.
